\section{ELF 目标文件、链接与装载}
\label{sec:elf}

\subsection{目标编码与重定位}

参考 RV64 汇编经汇编后形成 EM\_RISCV 类型的 ELF 可重定位对象。对象文件按节组织指令和元数据，符号表记录 \code{clamp\_input}、\code{factorial} 与 \code{main} 等已定义函数，以及 \code{getint}、\code{putint} 和 \code{putch} 三个未定义外部符号。ELF 通用规范定义节、符号和重定位的共同结构，RISC-V psABI 补充处理器相关的重定位类型\cite{elf-gabi-2025,riscv-psabi-1.0}。

对象偏移 \code{0x2a} 处对 \code{getint} 的调用具有 \code{R\_RISCV\_CALL\_PLT} 与配套 \code{R\_RISCV\_RELAX} 记录。前者关联调用位置和目标符号，后者允许链接器在最终布局已知后执行松弛（linker relaxation）。同样的结构也出现在其他外部调用处。可重定位对象已经包含目标指令字节，但跨对象调用的最终位移仍由链接阶段填写。

\subsection{符号解析、节布局与静态链接}

链接器解析输入对象与归档中的全局符号，选择所需成员，合并输入节并分配地址，随后应用重定位，生成带入口地址和程序头表的 ELF 可执行文件\cite{lld-design-2026,elf-gabi-2025}。SysY 运行时库在此提供三个输入输出符号的定义；启动代码提供 \code{\_start}，完成初始运行环境建立后调用 \code{main}。静态链接使程序执行所需代码进入同一文件。

表~\ref{tab:sizes} 给出四种入口表示产生的对象和静态可执行文件磁盘大小。C 与 SysY 对象相差 8 字节，参考 IR 和参考汇编对象依次更小；静态文件的差异仅为数百字节，因为共同的启动及输入输出代码占主要部分。文件大小同时受节、符号表、对齐和元数据影响，适用于产物结构比较。

\begin{table}[tb]
\caption{可重定位对象与静态可执行文件的磁盘大小（字节）。}
\label{tab:sizes}
\centering\small
\begin{tabular}{@{}lrr@{}}
\toprule
入口表示 & 可重定位对象 & 静态可执行文件\\
\midrule
C 观察形态 & 2,736 & 627,032\\
SysY 源 & 2,744 & 627,032\\
参考 LLVM IR & 2,312 & 626,968\\
参考 RV64 汇编 & 1,512 & 626,832\\
\bottomrule
\end{tabular}
\end{table}

链接松弛利用全局布局信息改写先前的通用地址构造或调用序列。因此，对象反汇编与最终文件反汇编可能具有不同编码，而其调用语义保持不变。

\subsection{程序头与进程映像}

Linux \code{execve} 根据 ELF 程序头表建立新进程映像。\code{PT\_LOAD} 项规定文件范围、虚拟地址、内存大小、对齐和访问权限；当 \code{p\_memsz} 大于 \code{p\_filesz} 时，额外内存区域以零初始化\cite{linux-elf-2026,linux-execve-2026}。内核完成映射和初始栈准备后，将控制转移至 ELF 入口 \code{\_start}。

\begin{table}[tb]
\caption{ELF 节视图与装载段视图的功能对应。}
\label{tab:sections-segments}
\centering\small
\begin{tabularx}{0.9\linewidth}{@{}lXl@{}}
\toprule
链接节 & 主要内容 & 典型装载段\\
\midrule
\code{.text}, \code{.rodata} & 指令与只读数据 & R--/R-X \code{PT\_LOAD}\\
\code{.data}, \code{.bss} & 初始化数据与零填充区域 & RW- \code{PT\_LOAD}\\
\code{.symtab}, \code{.strtab} & 链接和诊断元数据 & 通常不映射\\
\bottomrule
\end{tabularx}
\end{table}

节用于链接工具对文件内容分类与合并，段用于装载器建立内存映射。至此，源级函数和变量已转换为具有确定地址、权限和入口状态的机器程序。
